%% ch10_conclusions.tex
%% Chapter 10: Conclusions and Future Work
%% ============================================================
\chapter{Conclusions and Future Work}
\label{ch:conclusions}

\section{Summary of Contributions}
\label{sec:ch10_summary}

This thesis has developed, analysed, and validated a hierarchical
self-adaptive protection and control framework—the \textbf{Self-Adaptive
Multi-Bus Protection (SAMBP)} system—for power distribution and
sub-transmission networks dominated by Inverter-Based Resources.
The framework addresses the twelve failure modes of legacy relay
algorithms identified in Chapter~\ref{ch:challenges} through six chapters
of original technical contributions.

\subsection{Tier 1: Local Relay Physics}

\paragraph{Chapter~\ref{ch:line_diff} — Line Differential Protection.}
The physics-constrained adaptive 87L threshold (C1) derives the optimal
bias slope $k_m^\ast(\kibr)$ from the simultaneous stability and sensitivity
constraints, eliminating the need for conservative fixed settings.
Proposition~\ref{prop:ch4_87ln} proves that IBR negative-sequence
suppression paradoxically \emph{improves} the 87LN element's
sensitivity ($\rho_2 = 2$ independent of $\kibr$). The four 87L
elements (C1--C4) provide complementary coverage across all fault types.

\paragraph{Chapter~\ref{ch:distance} — Distance Protection.}
The IBR-corrected quadrilateral characteristic (C5) eliminates overreach
caused by reactive current injection by subtracting the IBR impedance
shift using the digital twin's $\alpha$ estimate. The Monte Carlo
framework (C6) provides the first probabilistic reliability quantification
for IBR distance protection across the full operating envelope.

\paragraph{Chapter~\ref{ch:microgrid} — Microgrid Protection.}
The GFL/GFM dual-mode characterisation (C7) quantifies the 4--8$\times$
fault current reduction at grid-to-island transition. The ROCOF islanding
detection scheme (C8) achieves $< 1\%$ NDZ and 114~ms isolation time.
The adaptive OC setting-group (C9) switches within 5~ms of GOOSE receipt.

\subsection{Tier 2: Digital Monitoring and Learning}

\paragraph{Chapter~\ref{ch:digital_twin} — Digital Twin.}
The non-blocking DT architecture (C10) decouples EKF computation from
relay logic using a shadow register, guaranteeing that DT updates
never delay primary trip decisions. The EKF joint estimation (C11)
achieves RMSE $\leq 0.024$~pu for $\kibr$, $\leq 0.5\%$ for $\Zline$,
and $\leq 1.1\%$ for fault location, bounded by the derived CRLB.

\paragraph{Chapter~\ref{ch:ml} — Physics-Informed ML.}
The PINN (C12) achieves 96.1\% 5-class fault classification accuracy
with zero physics-constraint violations by embedding three constraints
in the training loss function. Permutation importance (C13) shows that
physics-derived composite features are redundant; the minimum effective
feature set is $\{\kibr, \alpha_\mathrm{margin}\}$ (2 features, 94\%
accuracy). CUSUM-SGD online learning (C14) adapts to fleet composition
drift within 21 detection samples and 50 update steps with $< 2$~pp
catastrophic forgetting.

\subsection{Tier 3: Centralised Coordination}

\paragraph{Chapter~\ref{ch:centralised} — Centralised Architecture.}
The WAPC PMU fault localisation (C15) achieves 94.5\% correct-line
identification at $\sigma_V = 0.01$~pu using voltage-dip triangulation.
The CTI-graded settings engine (C16) produces 100\% CTI-compliant settings
in $< 3$~s. HMAC-SHA256 (C17) reduces GOOSE attack risk by 54.9\% at
0.72\% latency overhead. The closed-loop CUSUM$\to$settings$\to$GOOSE
protocol (C18) completes the full adaptive cycle in $\approx 31$~minutes
with 8/9 steps automated.

\section{Mapping to Research Objectives}
\label{sec:ch10_objectives}

Table~\ref{tab:ch10_objectives} maps each objective (Section~\ref{sec:ch1_objectives})
to the contribution(s) that address it, the method, and the key metric.

\begin{table}[htbp]
\caption{Research objectives achieved}
\label{tab:ch10_objectives}
\begin{tabularx}{\textwidth}{@{}lp{2cm}p{3.5cm}p{4cm}X@{}}
\toprule
Obj. & Chapter & Contribution & Key metric & Status \\
\midrule
O1 & \ref{ch:line_diff} & C1: Adaptive 87L
  & 98\% sensitivity, 0 false trips & Achieved \\
O2 & \ref{ch:distance} & C5: Quadrilateral 21
  & $P_\mathrm{detect} = 99.3\%$, $P_\mathrm{secure} = 99.97\%$ & Achieved \\
O3 & \ref{ch:microgrid} & C7--C9: Microgrid
  & 114~ms mode switch, NDZ $< 1\%$ & Achieved \\
O4 & \ref{ch:digital_twin} & C10--C11: EKF DT
  & RMSE: 0.024, 0.5\%, 1.1\% & Achieved \\
O5 & \ref{ch:ml} & C12--C14: PINN + CUSUM
  & 96.1\% accuracy, $< 2$~pp forgetting & Achieved \\
O6 & \ref{ch:centralised} & C15--C18: WAPC
  & 94.5\% localise, 100\% CTI, 54.9\% risk reduction & Achieved \\
\bottomrule
\end{tabularx}
\end{table}

\section{Tier-1 Contributions}
\label{sec:ch10_tier1}

Six contributions achieved Tier-1 status (novelty $= 5$, significance $= 5$):

\begin{description}
\item[C1 — Physics-Constrained Adaptive 87L Threshold:]
  The first relay algorithm that modulates the bias slope as a function
  of real-time IBR current ratio, with closed-form derivation from
  simultaneous stability and sensitivity constraints.

\item[C7 — GFL/GFM Dual-Mode Characterisation:]
  The first systematic characterisation of protection-relevant differences
  between GFL and GFM inverter operating modes, quantifying the 4--8$\times$
  fault current reduction and its implications for relay setting groups.

\item[C10 — Non-Blocking Digital Twin Architecture:]
  The first digital twin architecture for protection that guarantees
  DT updates never delay primary trip decisions through the shadow
  register non-blocking pattern.

\item[C12 — PINN Fault Classifier:]
  The first Physics-Informed Neural Network applied to power system
  protection, embedding three circuit-law constraints directly in the
  loss function and guaranteeing zero physics-constraint violations.

\item[C15 — Wide-Area PMU Fault Localisation:]
  The first fault localisation criterion based on maximising combined
  endpoint voltage dip, with Proposition~\ref{prop:ch9_localise} proving
  its correctness, combined with N-1 security verification.

\item[C18 — Closed-Loop Adaptive Protocol:]
  The first closed-loop adaptive protection protocol that chains
  drift detection, settings re-optimisation, and GOOSE re-authentication
  into a single automated workflow.
\end{description}

\section{Limitations}
\label{sec:ch10_limitations}

The following limitations apply to the work presented in this thesis:

\begin{enumerate}
\item \textbf{Single-substation scope:} The WAPC is designed for a single
      substation up to 34 buses. Multi-substation coordination (wide-area
      hierarchical protection) is not addressed.

\item \textbf{Synthetic training data:} All PINN training and validation
      data are synthetic, generated by the SAMBP simulation platform.
      The model has not been validated on real-field fault recordings.

\item \textbf{SIL validation only:} Physical HIL testing on commercial
      IEDs (SEL-411L) is pending hardware availability.

\item \textbf{Fixed forgetting factor:} The online learning uses a fixed
      $\lambda = 0.99$; an adaptive $\lambda$ that responds to drift
      velocity would improve convergence speed.

\item \textbf{VSM model simplification:} GFM inverters are modelled as
      VSMs; full virtual synchronous machine dynamics including AVR
      are not captured beyond $t = 150$~ms.
\end{enumerate}

\section{Future Work and Open Gaps}
\label{sec:ch10_gaps}

Seven research gaps are identified, of which two are high-priority:

\paragraph{G4 (High priority) — Multi-Substation WAPC Scaling.}
The current WAPC architecture is centralised at a single substation.
Extending to a wide-area network of 50--200 buses requires a hierarchical
architecture: local substation coordinators (analogous to the current
WAPC) reporting to a regional coordinator. The key technical challenge
is consensus under communication latency: how to agree on fault
location and settings updates when the regional coordinator receives
information from multiple substations with 10--50~ms propagation delays.

\paragraph{G7 (High priority) — Real Field Fault Data for PINN Validation.}
The synthetic training data may not capture the full diversity of
real IBR fault waveforms, including:
\begin{itemize}
\item Manufacturer-specific current limiter implementations.
\item Legacy IBRs not compliant with IEEE~2800.
\item High-impedance faults on resistive soil types.
\item Evolving faults (initial HIF developing into bolted fault).
\end{itemize}
A collaboration with a DSO operating an IBR-dominated network to collect
and label real fault recordings is the critical next step.

\paragraph{G1 (Medium) — Physical HIL on Commercial IEDs.}
The SIL jitter model is a proxy for physical relay behaviour. Testing
on SEL-411L or similar commercial IEDs would validate the SAMBP
algorithms against real relay firmware, including the undocumented
preprocessing and filtering that commercial relays apply.

\paragraph{G2 (Medium) — Variable Forgetting Factor.}
An adaptive $\lambda(t)$ that increases when CUSUM detects rapid drift
and decreases when the operating point is stable would improve both
adaptation speed and long-term accuracy stability.

\paragraph{G3 (Low) — \texorpdfstring{$Z_\mathrm{line}$}{Z\_line} Full Observability.}
The EKF cannot estimate $\Zline$ at fault location $\alpha = 0$ from
single-ended measurements (Proposition~\ref{prop:ch7_observability}).
Two-ended voltage phasors from PMUs would remove this observability gap
and improve the accuracy of the distance relay reach correction.

\paragraph{G5 (Medium) — GOOSE Flooding Mitigation.}
HMAC-SHA256 addresses spoofing and replay but not flooding (DoS).
Switch-level rate limiting and VLAN segmentation are the appropriate
countermeasures; these require changes to the substation LAN architecture
and must be validated against IEC~61850 real-time performance requirements.

\paragraph{G6 (Low) — Full VSM Dynamics.}
The GFM model used in this thesis is a simplified VSM; full virtual
synchronous generator dynamics including AVR excitation, turbine governor,
and power system stabiliser (PSS) become significant for island voltage
stability analysis beyond 150~ms. Including these dynamics would extend
the DT validity horizon and improve islanding detection accuracy.

\section{Concluding Remarks}
\label{sec:ch10_concluding}

The displacement of synchronous machines by IBRs is irreversible and
accelerating. The protection crisis it creates—missed trips, sympathetic
trips, false trips, and coordination collapses—is not a temporary
transition problem but a permanent structural feature of all future
high-IBR power systems. The SAMBP framework proposed in this thesis
addresses this crisis at all three levels at which it manifests:
in the relay physics (Tier~1), in the system state knowledge
(Tier~2), and in the network coordination (Tier~3).

The 18 original contributions, of which 6 are Tier-1, advance the
state of the art in each area. More importantly, their integration
into the closed-loop CUSUM$\to$settings$\to$GOOSE protocol demonstrates
that adaptive protection—protection that maintains its performance
as the fleet composition evolves—is technically achievable within
the constraints of existing IEC~61850 infrastructure and commercial
relay hardware.

The remaining open gaps (principally multi-substation scaling and
field data validation) define a clear post-doctoral research agenda
that this thesis has equipped with the foundational methodology,
validated algorithms, and documented performance bounds needed to
make rapid progress.
